% GENERATED FILE. Edit canonical lessons, then run npm run generate:books.
% canonical-source-hash: fnv1a64:d71de47877467688
% canonical-lessons: JA-C380-ekiben, JA-C380-shashou, JA-C380-nesugosu, JA-C380-utouto, JA-C380-hitokuchi

\chapter{Lunch on the Train}
\label{ch:ja-lunch-on-the-train}
\hlchaptermodality{380}

\begin{chapteropening}
\textbf{By the end of this chapter:} \emph{I can talk about eating and dozing on a long train ride.}

Complete the last lesson of chapter 380: I can talk about eating and dozing on a long train ride.
\end{chapteropening}

\section[\texorpdfstring{\ja{えきべん} (ekiben)}{えきべん (ekiben)}]{\ja{えきべん} (ekiben) — a station lunch box}

\label{lesson:JA-C380-ekiben}



\begin{quote}
Before the new one: say the Japanese for inside the train, then the Japanese for a fare.
\end{quote}

\subsection*{You'll want to know: \ja{えきべん}}
\textbf{\ja{えきべん}} — \emph{ekiben} — \textquotedblleft{}a station lunch box\textquotedblright{}.

\begin{quote}
\textbf{\ja{でんしゃで} \ja{えきべんを} \ja{たべました。}} — \emph{densha de ekiben o tabemashita} — \textquotedblleft{}I ate a station lunch box on the train.\textquotedblright{}
\end{quote}

\begin{cousinweb}
\textbf{\ja{えきべん}} (\emph{ekiben}) is Sino-Japanese (kango): two Chinese characters read with their on'yomi, readings borrowed from Chinese. They are \emph{eki}, \textquotedblleft{}station\textquotedblright{} (Middle Chinese \emph{yek}), and \emph{ben}, the first sign of \emph{bentō}, \textquotedblleft{}a boxed lunch\textquotedblright{}, a word whose older history is not settled (JMdict, KANJIDIC2; Middle Chinese in Baxter's transcription).
\end{cousinweb}

\begin{grammarlens}[title={What you've built}]
One more word for getting around.
\end{grammarlens}

\subsection*{Your turn}
\begin{itemize}
\raggedright
  \item \emph{Say it:} \emph{ekiben}
  \item \emph{Say it:} \emph{ekiben}, once more
  \item \emph{Say it:} \emph{densha de ekiben o tabemashita} — I ate a station lunch box on the train.
  \item \emph{Recall:} say the Japanese for inside the train, then the Japanese for a fare, then say \emph{ekiben} again
\end{itemize}

\begin{tcolorbox}[breakable,colback=teal!4,colframe=teal!35!black,title={Before you move on}]
What does \ja{えきべん} mean? (\textquotedblleft{}A station lunch box\textquotedblright{}.) Say it once more.
\end{tcolorbox}
\section[\texorpdfstring{\ja{しゃしょう} (shashō)}{しゃしょう (shashō)}]{\ja{しゃしょう} (shashō) — a train conductor}

\label{lesson:JA-C380-shashou}



\begin{quote}
Before the new one: say the Japanese for a fare, then the Japanese for a station lunch box.
\end{quote}

\subsection*{You'll want to know: \ja{しゃしょう}}
\textbf{\ja{しゃしょう}} — \emph{shashō} — \textquotedblleft{}a train conductor\textquotedblright{}.

\begin{quote}
\textbf{\ja{しゃしょうに} \ja{きっぷを} \ja{わたしました。}} — \emph{shashō ni kippu o watashimashita} — \textquotedblleft{}I handed my ticket to the conductor.\textquotedblright{}
\end{quote}

\begin{cousinweb}
\textbf{\ja{しゃしょう}} (\emph{shashō}) is Sino-Japanese (kango): two Chinese characters read with their on'yomi, readings borrowed from Chinese. They are \emph{sha}, \textquotedblleft{}car\textquotedblright{} (Middle Chinese \emph{tsyhae}) and \emph{shō}, \textquotedblleft{}administer\textquotedblright{} (\emph{tsyangX}) (KANJIDIC2; Middle Chinese in Baxter's transcription).
\end{cousinweb}

\begin{grammarlens}[title={What you've built}]
One more word for getting around.
\end{grammarlens}

\subsection*{Your turn}
\begin{itemize}
\raggedright
  \item \emph{Say it:} \emph{shashō}
  \item \emph{Say it:} \emph{shashō}, once more
  \item \emph{Say it:} \emph{shashō ni kippu o watashimashita} — I handed my ticket to the conductor.
  \item \emph{Recall:} say the Japanese for a fare, then the Japanese for a station lunch box, then say \emph{shashō} again
\end{itemize}

\begin{tcolorbox}[breakable,colback=teal!4,colframe=teal!35!black,title={Before you move on}]
What does \ja{しゃしょう} mean? (\textquotedblleft{}A train conductor\textquotedblright{}.) Say it once more.
\end{tcolorbox}
\section[\texorpdfstring{\ja{ねすごす} (nesugosu)}{ねすごす (nesugosu)}]{\ja{ねすごす} (nesugosu) — to oversleep; to sleep past one's stop}

\label{lesson:JA-C380-nesugosu}



\begin{quote}
Before the new one: say the Japanese for a station lunch box, then the Japanese for a train conductor.
\end{quote}

\subsection*{You'll want to know: \ja{ねすごす}}
\textbf{\ja{ねすごす}} — \emph{nesugosu} — \textquotedblleft{}to oversleep; to sleep past one's stop\textquotedblright{}.

\begin{quote}
\textbf{\ja{でんしゃで} \ja{ねすごして、しゅうてんまで} \ja{いきました。}} — \emph{densha de nesugoshite, shūten made ikimashita} — \textquotedblleft{}I slept past my stop and went all the way to the end of the line.\textquotedblright{}
\end{quote}

Polite \textbf{\ja{ねすごします}}, \emph{nesugoshimasu}; a godan verb. Past (polite): \textbf{\ja{あにが} \ja{でんしゃで} \ja{ねすごしました。}}, \emph{ani ga densha de nesugoshimashita}.

\begin{cousinweb}
\textbf{\ja{ねすごす}} (\emph{nesugosu}) is a native Japanese word (wago), inherited rather than borrowed from Chinese, built from \emph{ne}, \textquotedblleft{}sleeping\textquotedblright{} (from \emph{neru}) and \emph{sugosu}, \textquotedblleft{}to let pass\textquotedblright{}; its kanji were attached later as kun'yomi (JMdict).
\end{cousinweb}

\begin{grammarlens}[title={What you've built}]
One more word for getting around.
\end{grammarlens}

\subsection*{Your turn}
\begin{itemize}
\raggedright
  \item \emph{Say it:} \emph{nesugosu}
  \item \emph{Say it:} \emph{nesugosu}, once more
  \item \emph{Say it:} \emph{densha de nesugoshite, shūten made ikimashita} — I slept past my stop and went all the way to the end of the line.
  \item \emph{Recall:} say the Japanese for a station lunch box, then the Japanese for a train conductor, then say \emph{nesugosu} again
\end{itemize}

\begin{tcolorbox}[breakable,colback=teal!4,colframe=teal!35!black,title={Before you move on}]
What does \ja{ねすごす} mean? (\textquotedblleft{}To oversleep; to sleep past one's stop\textquotedblright{}.) Say it once more.
\end{tcolorbox}
\section[\texorpdfstring{\ja{うとうと} (utōto)}{うとうと (utōto)}]{\ja{うとうと} (utōto) — dozing off}

\label{lesson:JA-C380-utouto}



\begin{quote}
Before the new one: say the Japanese for a train conductor, then the Japanese for to oversleep.
\end{quote}

\subsection*{You'll want to know: \ja{うとうと}}
\textbf{\ja{うとうと}} — \emph{utōto} — \textquotedblleft{}dozing off\textquotedblright{}.

\begin{quote}
\textbf{\ja{でんしゃで} \ja{うとうとしました。}} — \emph{densha de utōto shimashita} — \textquotedblleft{}I dozed off on the train.\textquotedblright{}
\end{quote}

\begin{cousinweb}
\textbf{\ja{うとうと}} (\emph{utōto}) is a native mimetic word that pictures someone nodding off into a light sleep: an imitative Japanese form, inherited rather than borrowed, whose older origin is not settled (JMdict).
\end{cousinweb}

\begin{grammarlens}[title={What you've built}]
One more word for getting around.
\end{grammarlens}

\subsection*{Your turn}
\begin{itemize}
\raggedright
  \item \emph{Say it:} \emph{utōto}
  \item \emph{Say it:} \emph{utōto}, once more
  \item \emph{Say it:} \emph{densha de utōto shimashita} — I dozed off on the train.
  \item \emph{Recall:} say the Japanese for a train conductor, then the Japanese for to oversleep, then say \emph{utōto} again
\end{itemize}

\begin{tcolorbox}[breakable,colback=teal!4,colframe=teal!35!black,title={Before you move on}]
What does \ja{うとうと} mean? (\textquotedblleft{}Dozing off\textquotedblright{}.) Say it once more.
\end{tcolorbox}
\section[\texorpdfstring{\ja{ひとくち} (hitokuchi)}{ひとくち (hitokuchi)}]{\ja{ひとくち} (hitokuchi) — a mouthful; a bite}

\label{lesson:JA-C380-hitokuchi}



\begin{quote}
Before the new one: say the Japanese for to oversleep, then the Japanese for dozing off.
\end{quote}

\subsection*{You'll want to know: \ja{ひとくち}}
\textbf{\ja{ひとくち}} — \emph{hitokuchi} — \textquotedblleft{}a mouthful; a bite\textquotedblright{}.

\begin{quote}
\textbf{\ja{えきべんを} \ja{ひとくち} \ja{たべました。}} — \emph{ekiben o hitokuchi tabemashita} — \textquotedblleft{}I had a bite of the station lunch box.\textquotedblright{}
\end{quote}

\begin{cousinweb}
\textbf{\ja{ひとくち}} (\emph{hitokuchi}) is a native Japanese word (wago), inherited rather than borrowed from Chinese, built from \emph{hito}, \textquotedblleft{}one\textquotedblright{} and \emph{kuchi}, \textquotedblleft{}mouth\textquotedblright{}; its kanji were attached later as kun'yomi (JMdict).
\end{cousinweb}

\begin{grammarlens}[title={What you've built}]
One more word for getting around.
\end{grammarlens}

\subsection*{Your turn}
\begin{itemize}
\raggedright
  \item \emph{Say it:} \emph{hitokuchi}
  \item \emph{Say it:} \emph{hitokuchi}, once more
  \item \emph{Say it:} \emph{ekiben o hitokuchi tabemashita} — I had a bite of the station lunch box.
  \item \emph{Recall:} say the Japanese for to oversleep, then the Japanese for dozing off, then say \emph{hitokuchi} again
\end{itemize}

\begin{tcolorbox}[breakable,colback=teal!4,colframe=teal!35!black,title={Before you move on}]
What does \ja{ひとくち} mean? (\textquotedblleft{}A mouthful; a bite\textquotedblright{}.) Say it once more.
\end{tcolorbox}
